% One-time imported draft. Editable native source; no Markdown regeneration.
\section{Cross the signed-word
boundary}\label{004-bit-inspector-x-addition-boundary-cross-the-signed-word-boundary}

{\def\LTcaptype{table} % do not increment counter
\begin{longtable}[]{@{}
  >{\RaggedRight\arraybackslash}p{(\linewidth - 4\tabcolsep) * \real{0.3333}}
  >{\RaggedRight\arraybackslash}p{(\linewidth - 4\tabcolsep) * \real{0.3333}}
  >{\RaggedRight\arraybackslash}p{(\linewidth - 4\tabcolsep) * \real{0.3333}}@{}}
\toprule\noalign{}
\begin{minipage}[b]{\linewidth}\RaggedRight
Original signed value
\end{minipage} & \begin{minipage}[b]{\linewidth}\RaggedRight
Word-plus-one result
\end{minipage} & \begin{minipage}[b]{\linewidth}\RaggedRight
Doubleword-plus-one result
\end{minipage} \\
\midrule\noalign{}
\endhead
\bottomrule\noalign{}
\endlastfoot
-1 & 0 & 0 \\
2147483646 & 2147483647 & 2147483647 \\
2147483647 & -2147483648 & 2147483648 \\
\end{longtable}
}

At the final row, the full word result is ffffffff80000000 and the full
doubleword result is 0000000080000000. Both have the same low word.

The doubleword column starts from the signed load. Starting from the
unsigned load would be a different calculation: 4294967295 plus one is
4294967296 for the all-one word. Its low word is zero, so a low-word
check would miss that error.
